\section{论文概览}

\subsection{核心定位}

HY-Embodied-0.5 面向真实世界智能体，目标是构建一个能够处理视觉、语言、空间理解、具身推理和行动规划的 embodied foundation model。与普通视觉语言模型相比，它不只需要“看图并回答”，还要更强地理解物体位置、空间关系、动作轨迹、环境状态和任务步骤。

从现有笔记看，论文的主线可以概括为三层：第一，构建适合边缘部署的视觉前端 HY-ViT 2.0；第二，在多模态主干中通过 MoT 解耦视觉 token 与文本 token 的计算路径；第三，通过预训练、中期训练、SFT、RL、RFT 和 OPD 等阶段，把基础视觉语言能力进一步转化为具身推理与规划能力。

\subsection{一句话总结}

HY-Embodied-0.5 的核心思想是：用高效视觉骨干获得语言可读且细节保留的视觉表示，再用 modality-specific 的 MoT 主干减少视觉训练对语言能力的破坏，最后通过多阶段后训练把模型推向真实具身任务所需的推理、定位、轨迹和规划能力。

\subsection{本文关注的关键问题}

\begin{itemize}
    \item 视觉端如何既轻量又强。论文使用 400M 规模的 HY-ViT 2.0，并通过更强内部 ViT 蒸馏、tiny LLM 语言监督和视觉重建监督来提升表示质量。
    \item 视觉表示如何不丢细节。论文引入离散视觉表示、codebook 和 next-code prediction，让普通视觉 token 接受 patch-level 监督，而不是只依赖最终文本答案的间接梯度。
    \item 多模态主干如何避免语言能力退化。论文采用 Mixture-of-Transformers，使视觉 token 和文本 token 使用不同的 QKV/FFN 参数，从而缓解大量视觉训练对原始 LLM 语言能力的冲击。
    \item 训练策略如何服务具身任务。预训练阶段联合优化语言、局部视觉和全局视觉目标；后续阶段转向标准自回归语言损失，并通过 SFT、RL、RFT 和 OPD 强化复杂推理与能力迁移。
\end{itemize}

\subsection{阅读时需要区分的两类内容}

原始笔记中多次提醒：论文有些地方只给出高层描述，没有展开代码级公式。因此整理时需要区分“论文明确陈述”和“基于 VLM/Transformer 常见实现的合理技术解释”。例如 tiny LLM 如何具体提供语言监督、离散视觉 tokenizer 如何实现、MoT 内部 token 路由如何写成代码，这些细节在论文中并非全部完整展开，理解时应避免把合理推断误当成论文原文。
